% GENERATED FILE. Edit canonical lessons, then run npm run generate:books.
% canonical-source-hash: fnv1a64:87792e814c6c5408
% canonical-lessons: FA-C494-iyadat-kardan, FA-C494-moraqebat-kardan, FA-C494-parastari-kardan, FA-C494-salamati, FA-C494-khub-shodan

\chapter{Caring for the Sick}
\label{ch:fa-caring-for-the-sick}
\hlchaptermodality{494}

\begin{chapteropening}
\textbf{By the end of this chapter:} \emph{I can talk about looking after someone who is ill.}

Complete the last lesson of chapter 494: I can talk about looking after someone who is ill.
\end{chapteropening}

\section[\texorpdfstring{\fa{عیادت} \fa{کردن} ('iyâdat kardan)}{عیادت کردن ('iyâdat kardan)}]{\fa{عیادت} \fa{کردن} ('iyâdat kardan) — to visit (a sick person)}

\label{lesson:FA-C494-iyadat-kardan}



\begin{quote}
Before the new one: say the Persian for to get vaccinated, then the Persian for a check-up.
\end{quote}

\subsection*{You'll want to know: \fa{عیادت} \fa{کردن}}
\textbf{\fa{عیادت} \fa{کردن}} (\emph{'iyâdat kardan}) — \textquotedblleft{}to visit (a sick person)\textquotedblright{}.

\begin{quote}
\fa{همکارانش} \fa{در} \fa{بیمارستان} \fa{از} \fa{او} \fa{عیادت} \fa{کردند}.
\end{quote}

\begin{quote}
\emph{hamkârânash dar bimârestân az u 'iyâdat kardand.}
\end{quote}

\begin{quote}
His colleagues visited him in hospital.
\end{quote}

\emph{az kasi 'iyâdat kardan}: said only of visiting the sick. \emph{'iyâdat} + \emph{kardan}.

\begin{cousinweb}
\emph{'iyâdat} is a loan, usually traced to Arabic \emph{ʿiyāda}, 'visiting the sick'. EtymDB has no entry for it. \emph{kardan} is inherited Persian.
\end{cousinweb}

\begin{grammarlens}[title={What you've built}]
One more word for telling, explaining and describing at B1.
\end{grammarlens}

\subsection*{Your turn}
\begin{itemize}
\raggedright
  \item \emph{Say it:} \emph{'iyâdat kardan}
  \item \emph{Say it:} \emph{'iyâdat kardan}, once more
  \item \emph{Say it:} a sentence of your own with \emph{'iyâdat kardan}
  \item \emph{Recall:} say the Persian for to get vaccinated, then the Persian for a check-up, then say \emph{'iyâdat kardan} again
\end{itemize}

\begin{tcolorbox}[breakable,colback=teal!4,colframe=teal!35!black,title={Before you move on}]
What does \fa{عیادت} \fa{کردن} mean? (\textquotedblleft{}To visit (a sick person)\textquotedblright{}.) Say it once more.
\end{tcolorbox}
\section[\texorpdfstring{\fa{مراقبت} \fa{کردن} (morâqebat kardan)}{مراقبت کردن (morâqebat kardan)}]{\fa{مراقبت} \fa{کردن} (morâqebat kardan) — to look after, to take care of}

\label{lesson:FA-C494-moraqebat-kardan}



\begin{quote}
Before the new one: say the Persian for a check-up, then the Persian for to visit.
\end{quote}

\subsection*{You'll want to know: \fa{مراقبت} \fa{کردن}}
\textbf{\fa{مراقبت} \fa{کردن}} (\emph{morâqebat kardan}) — \textquotedblleft{}to look after, to take care of\textquotedblright{}.

\begin{quote}
\fa{او} \fa{از} \fa{مادر} \fa{پیرش} \fa{مراقبت} \fa{میکند}.
\end{quote}

\begin{quote}
\emph{u az mâdar-e pirash morâqebat mikonad.}
\end{quote}

\begin{quote}
She looks after her elderly mother.
\end{quote}

\emph{az kasi morâqebat kardan}. \emph{morâqeb bâsh} — take care.

\begin{cousinweb}
\emph{morâqebat} is a loan, usually traced to Arabic \emph{murāqaba}, 'watching over'. EtymDB has no entry for it. \emph{kardan} is inherited Persian.
\end{cousinweb}

\begin{grammarlens}[title={What you've built}]
One more word for telling, explaining and describing at B1.
\end{grammarlens}

\subsection*{Your turn}
\begin{itemize}
\raggedright
  \item \emph{Say it:} \emph{morâqebat kardan}
  \item \emph{Say it:} \emph{morâqebat kardan}, once more
  \item \emph{Say it:} a sentence of your own with \emph{morâqebat kardan}
  \item \emph{Recall:} say the Persian for a check-up, then the Persian for to visit, then say \emph{morâqebat kardan} again
\end{itemize}

\begin{tcolorbox}[breakable,colback=teal!4,colframe=teal!35!black,title={Before you move on}]
What does \fa{مراقبت} \fa{کردن} mean? (\textquotedblleft{}To look after, to take care of\textquotedblright{}.) Say it once more.
\end{tcolorbox}
\section[\texorpdfstring{\fa{پرستاری} \fa{کردن} (parastâri kardan)}{پرستاری کردن (parastâri kardan)}]{\fa{پرستاری} \fa{کردن} (parastâri kardan) — to nurse, to care for}

\label{lesson:FA-C494-parastari-kardan}



\begin{quote}
Before the new one: say the Persian for to visit, then the Persian for to look after.
\end{quote}

\subsection*{You'll want to know: \fa{پرستاری} \fa{کردن}}
\textbf{\fa{پرستاری} \fa{کردن}} (\emph{parastâri kardan}) — \textquotedblleft{}to nurse, to care for\textquotedblright{}.

\begin{quote}
\fa{وقتی} \fa{بیمار} \fa{بودم،} \fa{خواهرم} \fa{از} \fa{من} \fa{پرستاری} \fa{کرد}.
\end{quote}

\begin{quote}
\emph{vaqti bimâr budam, khâharam az man parastâri kard.}
\end{quote}

\begin{quote}
When I was ill, my sister nursed me.
\end{quote}

\emph{parastâr} 'nurse' + \emph{-i} + \emph{kardan}. The person nursed takes \emph{az}.

\begin{cousinweb}
\emph{parastâri} is built on \emph{parastâr}, 'nurse', which is inherited: EtymDB traces it to a Middle Persian word for a maidservant (EtymDB). The ending \emph{-i} and \emph{kardan} are native too.
\end{cousinweb}

\begin{grammarlens}[title={What you've built}]
One more word for telling, explaining and describing at B1.
\end{grammarlens}

\subsection*{Your turn}
\begin{itemize}
\raggedright
  \item \emph{Say it:} \emph{parastâri kardan}
  \item \emph{Say it:} \emph{parastâri kardan}, once more
  \item \emph{Say it:} a sentence of your own with \emph{parastâri kardan}
  \item \emph{Recall:} say the Persian for to visit, then the Persian for to look after, then say \emph{parastâri kardan} again
\end{itemize}

\begin{tcolorbox}[breakable,colback=teal!4,colframe=teal!35!black,title={Before you move on}]
What does \fa{پرستاری} \fa{کردن} mean? (\textquotedblleft{}To nurse, to care for\textquotedblright{}.) Say it once more.
\end{tcolorbox}
\section[\texorpdfstring{\fa{سلامتی} (salâmati)}{سلامتی (salâmati)}]{\fa{سلامتی} (salâmati) — health}

\label{lesson:FA-C494-salamati}



\begin{quote}
Before the new one: say the Persian for to look after, then the Persian for to nurse.
\end{quote}

\subsection*{You'll want to know: \fa{سلامتی}}
\textbf{\fa{سلامتی}} (\emph{salâmati}) — \textquotedblleft{}health\textquotedblright{}.

\begin{quote}
\fa{ورزش} \fa{برای} \fa{سلامتی} \fa{خوب} \fa{است}.
\end{quote}

\begin{quote}
\emph{varzesh barâ-ye salâmati khub ast.}
\end{quote}

\begin{quote}
Exercise is good for your health.
\end{quote}

Also a toast and a wish: \emph{be salâmati} 'cheers, to your health'.

\begin{cousinweb}
\emph{salâmati} is built in Persian on \emph{salâmat}, 'soundness', a loan usually traced to Arabic \emph{salāma}, from the same root s-l-m as \emph{salâm}, with the native ending \emph{-i}.
\end{cousinweb}

\begin{grammarlens}[title={What you've built}]
One more word for telling, explaining and describing at B1.
\end{grammarlens}

\subsection*{Your turn}
\begin{itemize}
\raggedright
  \item \emph{Say it:} \emph{salâmati}
  \item \emph{Say it:} \emph{salâmati}, once more
  \item \emph{Say it:} a sentence of your own with \emph{salâmati}
  \item \emph{Recall:} say the Persian for to look after, then the Persian for to nurse, then say \emph{salâmati} again
\end{itemize}

\begin{tcolorbox}[breakable,colback=teal!4,colframe=teal!35!black,title={Before you move on}]
What does \fa{سلامتی} mean? (\textquotedblleft{}Health\textquotedblright{}.) Say it once more.
\end{tcolorbox}
\section[\texorpdfstring{\fa{خوب} \fa{شدن} (khub shodan)}{خوب شدن (khub shodan)}]{\fa{خوب} \fa{شدن} (khub shodan) — to get better, to recover}

\label{lesson:FA-C494-khub-shodan}



\begin{quote}
Before the new one: say the Persian for to nurse, then the Persian for health.
\end{quote}

\subsection*{You'll want to know: \fa{خوب} \fa{شدن}}
\textbf{\fa{خوب} \fa{شدن}} (\emph{khub shodan}) — \textquotedblleft{}to get better, to recover\textquotedblright{}.

\begin{quote}
\fa{او} \fa{بعد} \fa{از} \fa{یک} \fa{هفته} \fa{خوب} \fa{شد}.
\end{quote}

\begin{quote}
\emph{u ba'd az yek hafte khub shod.}
\end{quote}

\begin{quote}
He got better after a week.
\end{quote}

\emph{khub} 'good, well' + \emph{shodan} 'to become'. It also means 'to turn out well'.

\begin{cousinweb}
\emph{khub shodan} is built on two inherited words: \emph{khub}, 'good', which EtymDB traces to Middle Persian \emph{xūb} (EtymDB), and \emph{shodan}, 'to become'.
\end{cousinweb}

\begin{grammarlens}[title={What you've built}]
One more word for telling, explaining and describing at B1.
\end{grammarlens}

\subsection*{Your turn}
\begin{itemize}
\raggedright
  \item \emph{Say it:} \emph{khub shodan}
  \item \emph{Say it:} \emph{khub shodan}, once more
  \item \emph{Say it:} a sentence of your own with \emph{khub shodan}
  \item \emph{Recall:} say the Persian for to nurse, then the Persian for health, then say \emph{khub shodan} again
\end{itemize}

\begin{tcolorbox}[breakable,colback=teal!4,colframe=teal!35!black,title={Before you move on}]
What does \fa{خوب} \fa{شدن} mean? (\textquotedblleft{}To get better, to recover\textquotedblright{}.) Say it once more.
\end{tcolorbox}
